\documentclass[10pt,a4paper]{amsart}
\usepackage{amsmath,amssymb,amsthm,graphicx}
\usepackage{verbatim}
\newenvironment{reference}{\comment}{\endcomment}
%\newenvironment{reference}{}{}
\newenvironment{slogan}{\comment}{\endcomment}
\newenvironment{history}{\comment}{\endcomment}

% For commutative diagrams we use Xy-pic
\usepackage[all]{xy}

% We use 2cell for 2-commutative diagrams.
\xyoption{2cell}
% Theorem environments.
%
\theoremstyle{plain}
\newtheorem{theorem}[subsection]{Theorem}
\newtheorem{proposition}[subsection]{Proposition}
\newtheorem{lemma}[subsection]{Lemma}

\theoremstyle{definition}
\newtheorem{definition}[subsection]{Definition}
\newtheorem{example}[subsection]{Example}
\newtheorem{exercise}[subsection]{Exercise}
\newtheorem{situation}[subsection]{Situation}

\theoremstyle{remark}
\newtheorem{remark}[subsection]{Remark}
\newtheorem{remarks}[subsection]{Remarks}

\numberwithin{equation}{subsection}

% Macros
%
\def\lim{\mathop{\mathrm{lim}}\nolimits}
\def\colim{\mathop{\mathrm{colim}}\nolimits}
\def\Spec{\mathop{\mathrm{Spec}}}
\def\Hom{\mathop{\mathrm{Hom}}\nolimits}
\def\Ext{\mathop{\mathrm{Ext}}\nolimits}
\def\SheafHom{\mathop{\mathcal{H}\!\mathit{om}}\nolimits}
\def\SheafExt{\mathop{\mathcal{E}\!\mathit{xt}}\nolimits}
\def\Sch{\mathit{Sch}}
\def\Mor{\mathop{\mathrm{Mor}}\nolimits}
\def\Ob{\mathop{\mathrm{Ob}}\nolimits}
\def\Sh{\mathop{\mathit{Sh}}\nolimits}
\def\NL{\mathop{N\!L}\nolimits}
\def\CH{\mathop{\mathrm{CH}}\nolimits}
\def\proetale{{pro\text{-}\acute{e}tale}}
\def\etale{{\acute{e}tale}}
\def\QCoh{\mathit{QCoh}}
\def\Ker{\mathop{\mathrm{Ker}}}
\def\Im{\mathop{\mathrm{Im}}}
\def\Coker{\mathop{\mathrm{Coker}}}
\def\Coim{\mathop{\mathrm{Coim}}}

% Boxtimes
%
\DeclareMathSymbol{\boxtimes}{\mathbin}{AMSa}{"02}

%
% Macros for moduli stacks/spaces
%
\def\QCohstack{\mathcal{QC}\!\mathit{oh}}
\def\Cohstack{\mathcal{C}\!\mathit{oh}}
\def\Spacesstack{\mathcal{S}\!\mathit{paces}}
\def\Quotfunctor{\mathrm{Quot}}
\def\Hilbfunctor{\mathrm{Hilb}}
\def\Curvesstack{\mathcal{C}\!\mathit{urves}}
\def\Polarizedstack{\mathcal{P}\!\mathit{olarized}}
\def\Complexesstack{\mathcal{C}\!\mathit{omplexes}}
% \Pic is the operator that assigns to X its picard group, usage \Pic(X)
% \Picardstack_{X/B} denotes the Picard stack of X over B
% \Picardfunctor_{X/B} denotes the Picard functor of X over B
\def\Pic{\mathop{\mathrm{Pic}}\nolimits}
\def\Picardstack{\mathcal{P}\!\mathit{ic}}
\def\Picardfunctor{\mathrm{Pic}}
\def\Deformationcategory{\mathcal{D}\!\mathit{ef}}
\usepackage{cleveref}
\crefname{originaldefinition}{original definition}{original definitions}
\Crefname{originaldefinition}{Original definition}{Original definitions}
\makeatletter
\newlabel{pione-definition-fundamental-group}{{0BNC}{source}{}{sourcepointer-1}{}}
\newlabel{pione-definition-fundamental-group@cref}{{[originaldefinition][0][]0BNC}{[1][0][]source}{}{}{}}
\makeatother
\crefname{originalsection}{original section}{original sections}
\Crefname{originalsection}{Original section}{Original sections}
\makeatletter
\newlabel{section-locally-constant}{{09Y8}{source}{}{sourcepointer-2}{}}
\newlabel{section-locally-constant@cref}{{[originalsection][0][]09Y8}{[1][0][]source}{}{}{}}
\makeatother
\crefname{originaldefinition}{original definition}{original definitions}
\Crefname{originaldefinition}{Original definition}{Original definitions}
\makeatletter
\newlabel{pione-definition-G-set-continuous}{{04JJ}{source}{}{sourcepointer-3}{}}
\newlabel{pione-definition-G-set-continuous@cref}{{[originaldefinition][0][]04JJ}{[1][0][]source}{}{}{}}
\makeatother
\crefname{originaldefinition}{original definition}{original definitions}
\Crefname{originaldefinition}{Original definition}{Original definitions}
\makeatletter
\newlabel{definition-G-module-continuous}{{04JP}{source}{}{sourcepointer-4}{}}
\newlabel{definition-G-module-continuous@cref}{{[originaldefinition][0][]04JP}{[1][0][]source}{}{}{}}
\makeatother
\crefname{originallemma}{original lemma}{original lemmas}
\Crefname{originallemma}{Original lemma}{Original lemmas}
\makeatletter
\newlabel{lemma-characterize-finite-locally-constant}{{03RV}{source}{}{sourcepointer-5}{}}
\newlabel{lemma-characterize-finite-locally-constant@cref}{{[originallemma][0][]03RV}{[1][0][]source}{}{}{}}
\makeatother
\crefname{originaltheorem}{original theorem}{original theorems}
\Crefname{originaltheorem}{Original theorem}{Original theorems}
\makeatletter
\newlabel{pione-theorem-fundamental-group}{{0BND}{source}{}{sourcepointer-6}{}}
\newlabel{pione-theorem-fundamental-group@cref}{{[originaltheorem][0][]0BND}{[1][0][]source}{}{}{}}
\makeatother
\title{Finite locally constant sheaves on a connected scheme are finite representations of the etale fundamental group}
\author{The Stacks Project Authors}
\date{stacks-project:etale-cohomology.tex:local-snapshot-20261011; GNU FDL 1.2 or later}
\begin{document}
\maketitle
\noindent\textit{Editorial note.} All delivered mathematical text is the author's own, copied from the pinned Stacks Project chapter file etale-cohomology.tex; no statement or proof step was rewritten, reordered or removed.\par\medskip
\noindent\textit{Editorial note.} The proof is the author's complete argument for the equivalence, with both directions and the verification that they are inverse inside the counted span.\par\medskip
\noindent\textit{Editorial note.} Every labelled reference occurring inside the counted spans is carried in assembly.external\_references as a non-load-bearing navigation pointer (same-chapter ones at their definition site, cross-chapter ones at the invocation inside the counted span), so no load-bearing mathematics is replaced by a stub. No bibliography citation occurs inside the counted spans.\par\medskip
\noindent\textit{Editorial note.} The checked chapter list for same-name collisions is the set of stacks chapters staged in this repository: chow, cohomology, divisors, etale-cohomology, flat, more-algebra, more-morphisms, morphisms, spaces-more-morphisms, varieties together with this card's chapter etale-cohomology.\par\medskip
\noindent\textit{Editorial note.} The preamble is the chapter's own native preamble excerpt already used for this source: the Xy-pic setup (native preamble lines 14-18) and the native theorem declarations with the mathematical macros (native preamble lines 154-217).\par\medskip
\noindent\textit{Allowed standard prerequisites (editorial).}\begin{itemize}
\item Etale fundamental group of a connected scheme and its action on the fibres of a locally constant sheaf. Same chapter, the definition of the etale fundamental group at a geometric point and the description of locally constant sheaves by their fibres with the action of the group. Supplies the two categories between which the equivalence is asserted.
\item Trivialisation of a locally constant sheaf by a finite etale covering. Same chapter, the description of a finite locally constant sheaf as one trivialised by a finite etale covering, which realises its representation through a finite quotient of the fundamental group. Supplies the construction of the functor from sheaves to representations and its inverse.
\end{itemize}\medskip
\ifdefined\hypertarget\hypertarget{sourcepointer-1}{}\fi
\noindent\textit{Source reference.} Original definition 0BNC: sources/stacks-etale-cohomology/pinned-original-source.tex, source lines 10173--10173. Navigation pointer invoked at proof line 10173; the counted span contains the pointer only and the cited result is not reproduced inside it, so the pointer is registered as a non-load-bearing navigation stub in the established form for Stacks material. No bibliography citation is involved.\par\medskip
\ifdefined\hypertarget\hypertarget{sourcepointer-2}{}\fi
\noindent\textit{Source reference.} Original section 09Y8: sources/stacks-etale-cohomology/pinned-original-source.tex, source lines 9909--9909. Navigation pointer invoked at proof line 10175; the counted span contains the pointer only and the cited result is not reproduced inside it, so the pointer is registered as a non-load-bearing navigation stub in the established form for Stacks material. No bibliography citation is involved.\par\medskip
\ifdefined\hypertarget\hypertarget{sourcepointer-3}{}\fi
\noindent\textit{Source reference.} Original definition 04JJ: sources/stacks-etale-cohomology/pinned-original-source.tex, source lines 10177--10177. Navigation pointer invoked at proof line 10177; the counted span contains the pointer only and the cited result is not reproduced inside it, so the pointer is registered as a non-load-bearing navigation stub in the established form for Stacks material. No bibliography citation is involved.\par\medskip
\ifdefined\hypertarget\hypertarget{sourcepointer-4}{}\fi
\noindent\textit{Source reference.} Original definition 04JP: sources/stacks-etale-cohomology/pinned-original-source.tex, source lines 8786--8786. Navigation pointer invoked at proof line 10179; the counted span contains the pointer only and the cited result is not reproduced inside it, so the pointer is registered as a non-load-bearing navigation stub in the established form for Stacks material. No bibliography citation is involved.\par\medskip
\ifdefined\hypertarget\hypertarget{sourcepointer-5}{}\fi
\noindent\textit{Source reference.} Original lemma 03RV: sources/stacks-etale-cohomology/pinned-original-source.tex, source lines 9994--9994. Navigation pointer invoked at proof line 10184; the counted span contains the pointer only and the cited result is not reproduced inside it, so the pointer is registered as a non-load-bearing navigation stub in the established form for Stacks material. No bibliography citation is involved.\par\medskip
\ifdefined\hypertarget\hypertarget{sourcepointer-6}{}\fi
\noindent\textit{Source reference.} Original theorem 0BND: sources/stacks-etale-cohomology/pinned-original-source.tex, source lines 10185--10185. Navigation pointer invoked at proof line 10185; the counted span contains the pointer only and the cited result is not reproduced inside it, so the pointer is registered as a non-load-bearing navigation stub in the established form for Stacks material. No bibliography citation is involved.\par\medskip
\begin{lemma}
\label{lemma-locally-constant-on-connected}
Let $X$ be a connected scheme. Let $\overline{x}$ be a geometric point of $X$.
\begin{enumerate}
\item There is an equivalence of categories
$$
\left\{
\begin{matrix}
\text{finite locally constant}\\
\text{sheaves of sets on }X_\etale
\end{matrix}
\right\}
\longleftrightarrow
\left\{
\begin{matrix}
\text{finite }\pi_1(X, \overline{x})\text{-sets}
\end{matrix}
\right\}
$$
\item There is an equivalence of categories
$$
\left\{
\begin{matrix}
\text{finite locally constant}\\
\text{sheaves of abelian groups on }X_\etale
\end{matrix}
\right\}
\longleftrightarrow
\left\{
\begin{matrix}
\text{finite }\pi_1(X, \overline{x})\text{-modules}
\end{matrix}
\right\}
$$
\item Let $\Lambda$ be a finite ring. There is an equivalence of categories
$$
\left\{
\begin{matrix}
\text{finite type, locally constant}\\
\text{sheaves of }\Lambda\text{-modules on }X_\etale
\end{matrix}
\right\}
\longleftrightarrow
\left\{
\begin{matrix}
\text{finite }\pi_1(X, \overline{x})\text{-modules endowed}\\
\text{with commuting }\Lambda\text{-module structure}
\end{matrix}
\right\}
$$
\end{enumerate}
\end{lemma}

\begin{proof}
We observe that $\pi_1(X, \overline{x})$ is a profinite
topological group, see Fundamental Groups, Definition
\ref{pione-definition-fundamental-group}.
The left hand categories are defined in
Section \ref{section-locally-constant}.
The notation used in the right hand categories is taken from
Fundamental Groups, Definition \ref{pione-definition-G-set-continuous}
for sets and
Definition \ref{definition-G-module-continuous} for abelian groups.
This explains the notation.

\medskip\noindent
Assertion (1) follows from
Lemma \ref{lemma-characterize-finite-locally-constant}
and Fundamental Groups, Theorem \ref{pione-theorem-fundamental-group}.
Parts (2) and (3) follow immediately from this by endowing the underlying
(sheaves of) sets with additional structure. For example, a finite
locally constant sheaf of abelian groups on $X_\etale$ is the same thing
as a finite locally constant sheaf of sets $\mathcal{F}$
together with a map $+ : \mathcal{F} \times \mathcal{F} \to \mathcal{F}$
satisfying the usual axioms. The equivalence in (1) sends products
to products and hence sends $+$ to an addition on the corresponding
finite $\pi_1(X, \overline{x})$-set. Since $\pi_1(X, \overline{x})$-modules
are the same thing as $\pi_1(X, \overline{x})$-sets with a compatible
abelian group structure we obtain (2). Part (3) is proved in
exactly the same way.
\end{proof}

\end{document}
